\chapter{Bone Density Scan (DEXA)}

\section*{What Is This Test?}
A bone density scan, also called dual-energy X-ray absorptiometry (DEXA or DXA), uses very low-dose X-rays to measure bone mineral density, usually at the hip and lumbar spine. It helps diagnose low bone mass and osteoporosis, estimate fracture risk, and monitor response to treatment.

\section*{Alternative Names}
Bone Mineral Density Test; BMD Test; DEXA Scan; DXA Scan; Bone Densitometry

\section*{Why It Is Ordered}
Assessment of fracture risk; suspected osteoporosis or osteopenia; a previous low-trauma fracture; long-term glucocorticoid use; conditions or medicines that weaken bone; and monitoring of osteoporosis therapy. Screening recommendations depend on age, sex, fracture risk, and local guidelines.

\section*{How This Test Is Performed}
The patient lies on a padded table while a scanning arm passes over the hip and spine. The patient remains still, and the scan usually takes 10--30 minutes. Results are reported using bone-density scores interpreted with age, sex, menopausal status, fracture history, and clinical risk factors.

\section*{Preparation, Risks, and Limitations}
Calcium supplements may need to be withheld for 24--48 hours, and metal objects should be removed as directed. Tell the provider about possible pregnancy and recent barium or contrast examinations. Radiation exposure is very low. Arthritis, spinal deformity, fractures, hardware, and severe body-size limitations can affect interpretation; the scan estimates risk but does not measure bone strength perfectly.

\section*{References}
\begin{enumerate}
  \item MedlinePlus. Bone Density Scan.
  \item International Society for Clinical Densitometry. Official Positions.
\end{enumerate}

\clearpage
